% ch14 译文：原文附录 A §A.b（main.tex 第 5024--5545 行）
% [原文 p.50]
\begin{center}
\resizebox{0.98\textwidth}{!}{$\begin{aligned}
&
\check{C}_{G}^p
\Big(
\mathfrak{U},
\mathcal{H}om_{\mathcal{O}_X}^q
\big(
(\mathcal{F}^{\bullet},\theta^{\bullet}),
(\mathcal{F}'^{\bullet},\theta'^{\bullet})
\big)
\Big)
\cong
\prod_{i_0<\cdots<i_p}
\bigoplus_{i\in\mathbb{Z}}
\Hom_{\VB_{G}(U_{i_0\cdots i_p})}
\big(
(\mathcal{F}^i,\theta^i)|_{U_{i_0\cdots i_p}}),
(\mathcal{F}'^i,\theta'^{i+q}|_{U_{i_0\cdots i_p}})
\big),
\\[2mm]
&
\check{C}_{H}^r
\Big(
\mathfrak{V},
\mathcal{H}om_{\mathcal{O}_Y}^s
\big(
(\mathcal{G}^{\bullet},\eta^{\bullet}),
(\mathcal{G}'^{\bullet},\eta'^{\bullet})
\big)
\Big)
\cong
\prod_{j_0<\cdots<j_r}
\bigoplus_{j\in\mathbb{Z}}
\Hom_{\VB_{H}(V_{j_0\cdots j_p})}
\big(
(\mathcal{G}^j,\eta^j)|_{V_{j_0\cdots j_r}}),
(\mathcal{G}'^j,\eta'^{j+s}|_{V_{j_0\cdots j_r}})
\big),
\\[2mm]
&
\check{C}_{G\times H}^t
\Big(
\mathfrak{U}\times\mathfrak{V},
\mathcal{H}om_{\mathcal{O}_{X\times Y}}^u
\big(
(\mathcal{F}^{\bullet},\theta^{\bullet})
\boxtimes
(\mathcal{G}^{\bullet},\eta^{\bullet}),
(\mathcal{F}'^{\bullet},\theta'^{\bullet})
\boxtimes
(\mathcal{G}'^{\bullet},\eta'^{\bullet})
\big)
\Big)
\\[2mm]
&
\hspace{-5mm}
\cong
\prod_{\alpha_0<\cdots<\alpha_t}
\bigoplus_{k\in\mathbb{Z}}
\bigoplus_{\substack{
i+j=k
\\[1mm]
q+s=u
}}
\Hom_{\VB_{G\times H}(W_{\alpha_0\cdots \alpha_t})}
\Big(
\big(
(\mathcal{F}^i,\theta^i)
\boxtimes
(\mathcal{G}^j,\eta^j)
\big)|_{W_{\alpha_0\cdots \alpha_t}},
\big(
(\mathcal{F}'^{i+q},\theta'^{i+q})
\boxtimes
(\mathcal{G}'^{j+s},\eta'^{j+s})
\big)|_{W_{\alpha_0\cdots \alpha_t}})
\Big)
\end{aligned}$}
\end{center}

\begin{yizhu}[等同式里的指标笔误]
第二个等同式里 $\Hom_{\VB_{H}(V_{j_0\cdots j_p})}$ 的下标应为 $V_{j_0\cdots j_r}$：这是
$Y$ 侧的 \v{C}ech 项，次数是 $r$，不应出现 $X$ 侧的 $p$。第一个等同式右端的
$(\mathcal{F}'^i,\theta'^{i+q}|_{U_{i_0\cdots i_p}})$ 也照原文排印，但应读成
$(\mathcal{F}'^{i+q},\theta'^{i+q})|_{U_{i_0\cdots i_p}}$，即 $q$ 是内层态射的次数、
$i$ 是对象次数。两处都是原稿的排版笔误，不影响后文推导；译文照原样保留，以免改动公式。
\end{yizhu}

我们把
\begin{center}
\resizebox{\textwidth}{!}{$\begin{aligned}
C_{\otimes}^{n}
&=
\bigoplus_{l+m=n}
\Big(
\hom_{\Perf_{G}(X;\mathfrak{U})_{\dg}}^l
\big(
(\mathcal{F}^{\bullet},\theta^{\bullet}),
(\mathcal{F}'^{\bullet},\theta'^{\bullet})
\big)
\underset{\field}{\otimes}
\hom_{\Perf_{H}(Y;\mathfrak{V})_{\dg}}^m
\big(
(\mathcal{G}^{\bullet},\eta^{\bullet}),
(\mathcal{G}'^{\bullet},\eta'^{\bullet})
\big)
\Big)
\\[2mm]
&
\hspace{-5mm}
\cong
\bigoplus_{l+m=n}
\Bigg(
\bigg(
\bigoplus_{p+q=l}
\check{C}_{G}^p
\Big(
\mathfrak{U},
\mathcal{H}om_{\mathcal{O}_X}^q
\big(
(\mathcal{F}^{\bullet},\theta^{\bullet}),
(\mathcal{F}'^{\bullet},\theta'^{\bullet})
\big)
\Big)
\bigg)
\underset{\field}{\otimes}
\bigg(
\bigoplus_{r+s=m}
\check{C}_{H}^r
\Big(
\mathfrak{V},
\mathcal{H}om_{\mathcal{O}_Y}^s
\big(
(\mathcal{G}^{\bullet},\eta^{\bullet}),
(\mathcal{G}'^{\bullet},\eta'^{\bullet})
\big)
\Big)
\bigg)
\Bigg)
\\[2mm]
&
\hspace{-5mm}
\cong
\bigoplus_{l+m=n}
\bigoplus_{\substack{
p+q=l
\\[1mm]
r+s=m
}}
\prod_{\substack{
i_0<\cdots<i_p
\\[0mm]
j_0<\cdots<j_r
}}
\bigoplus_{\substack{
i\in\mathbb{Z}
\\[1mm]
j\in\mathbb{Z}
}}
\Big(
\Hom_{\VB_{G}(U_{i_0\cdots i_p})}
\big(
(\mathcal{F}^i,\theta^i)|_{U_{i_0\cdots i_p}},
(\mathcal{F}'^{i+q},\theta'^{i+q})|_{U_{i_0\cdots i_p}}
\big)
\\
&\phantom{
\cong
\bigoplus_{l+m=n}
\bigoplus_{\substack{
p+q=l
\\[1mm]
r+s=m
}}
\prod_{\substack{
i_0<\cdots<i_p
\\[0mm]
j_0<\cdots<j_r
}}
\bigoplus_{\substack{
i\in\mathbb{Z}
\\[1mm]
j\in\mathbb{Z}
}}
\Big(
}
\underset{\field}{\otimes}
\Hom_{\VB_{H}(V_{j_0\cdots j_r})}
\big(
(\mathcal{G}^j,\eta^j)|_{V_{j_0\cdots j_r}},
(\mathcal{G}'^{j+s},\eta'^{j+s})|_{V_{j_0\cdots j_s}}
\big)
\Big)
\end{aligned}$}
\end{center}
中多重次数为 $(p,q,r,s)$ 的齐次元素记为
$$
f\otimes g
=
\Big(
f_{i_0\cdots i_p}^i
\otimes
g_{j_0\cdots j_r}^j
\Big)
_{\substack{
i,j\in\mathbb{Z}
\\
i_0<\cdots<i_p
\\
j_0<\cdots<j_r
}}.
$$

(\ref{eqn:EZMaps}) 中的 Eilenberg--Zilber 映射 $S$、$T$ 与 $H$ 诱导出
$C_{\otimes}^{\bullet}$ 与 $C_{\times}^{\bullet}$ 之间的上链映射 $S^*$、$T^*$ 与 $H^*$：
$$
\begin{tikzcd}[arrow style=tikz,>=stealth,row sep=2.5em,column sep=3.5em]
C_{\otimes}^{\bullet}
  \arrow[r, shift left, <-, "S^*"]
&
C_{\times}^{\bullet}
  \arrow[l, shift left, <-, "T^*"]
  \arrow[loop right, "\mathbb{H}^{\ast}"]
\end{tikzcd}
$$

\begin{yizhu}[方向警告：$S^*$ 与 $S$ 反着走]
容易弄反：链映射 $S$ 走 $C_{\bullet}(I)\otimes C_{\bullet}(J)\to C_{\bullet}(I\times J)$，而
$\Hom$ 是反变的，所以诱导出的 $S^{\ast}$ 反向走 $C_{\times}^{\bullet}\to C_{\otimes}^{\bullet}$；
$T$ 与 $T^{\ast}$ 同理反向，图里指向 $C_{\otimes}^{\bullet}$ 的箭头才写作 $S^*$。记住
$T^{\ast}S^{\ast}=\id$（作用在 $C_{\times}^{\bullet}$ 上）、$S^{\ast}T^{\ast}$ 只与恒等相差一个同伦，
就与链层面的 $S\comp T=\id$、$T\comp S\sim\id$ 完全对应。
\end{yizhu}

% [原文 p.51]
\begin{dfn}
(1) 链映射
$
T:
C_{\bullet}(I\times J)
\rightarrow
C_{\bullet}(I)\otimes C_{\bullet}(J)
$
诱导上链复形之间的映射
\begin{align*}
T^*:
C_{\otimes}^n
&\rightarrow
C_{\times}^n
\\
\quad
f\otimes g
=
\Big(
f_{i_0\cdots i_p}^i
\otimes
g_{j_0\cdots j_r}^j
\Big)
_{\substack{
i,j\in\mathbb{Z}
\\
i_0<\cdots<i_p
\\
j_0<\cdots<j_r
}}
&\mapsto
T^*(f\otimes g)
:=
\Big(
T^*(f\otimes g)_{\alpha_0\cdots\alpha_t}^{i,j}
\Big)
_{\substack{
i,j\in\mathbb{Z}
\\
i_0<\cdots<i_p
\\
j_0<\cdots<j_r
}},
\end{align*}
其分量由
$$
T^*(f\otimes g)_{\alpha_0\cdots\alpha_t}^{i,j}
:=
\sum_{\substack{
i_0<\cdots<i_p
\\
j_0<\cdots<j_r
}}
(-1)^{qr+qs+si}
T_{\alpha_0\cdots\alpha_t}^{i_0\cdots i_p;j_0\cdots j_r}
\big(
f_{i_0\cdots i_p}^i
\boxtimes
g_{j_0\cdots j_r}^j
\big)\big|_{W_{\alpha_0\cdots\alpha_t}},
$$
定义，其中系数
$
T_{\alpha_0\cdots\alpha_t}^{i_0\cdots i_p;j_0\cdots j_r}
\in\mathbb{Z}
$
由公式
$$
T\big((\alpha_0,\dots,\alpha_t)\big)
=
\sum_{\substack{
i_0<\cdots<i_p
\\
j_0<\cdots<j_r
}}
T_{\alpha_0\cdots\alpha_t}^{i_0\cdots i_p;j_0\cdots j_r}
(i_0,\dots,i_p)\otimes(j_0,\dots,j_r).
$$
确定。由 $T$ 的定义（定义 \ref{defn:EZMaps}.(2)）可知
$$
T_{\alpha_0\cdots\alpha_t}^{i_0\cdots i_p;j_0\cdots j_r}
=
\begin{cases}
1
&
\text{若 }
(\alpha_0,\dots,\alpha_t)
=
\big((i_0,*),\dots,(i_{p-1},*),(i_p,j_0),(*,j_1),\dots,(*,j_r)\big),
\\
0
&
\text{否则},
\end{cases}
$$
因此
$$
T^*(f\otimes g)_{\alpha_0\cdots\alpha_t}^{i,j}
=
(-1)^{qr+qs+si}
\big(
f_{i_0\cdots i_p}^i
\boxtimes
g_{j_p\cdots j_t}^j
\big)\big|_{W_{\alpha_0\cdots\alpha_t}}
$$
其中 $(\alpha_0,\dots,\alpha_t)=\big((i_0,j_0),\dots,(i_t,j_t)\big)$。

\begin{yizhu}[符号 $(-1)^{qr+qs+si}$ 怎么读]
这个符号不是随手写的：$T^{\ast}$ 要与两个复形的微分交换（命题 \ref{prp:STHd-relation} 第四条），
而两边的微分都按 Koszul 法则带符号，换序时便补出 $(-1)^{qr+qs+si}$。它可以读成两次分次交换：
$(-1)^{q(r+s)}$ 来自内层次数 $q$ 越过 $g$ 的整体次数 $r+s$，$(-1)^{si}$ 来自对象次数 $i$ 越过
$g$ 的内层次数 $s$。逐项验算一次即可确认；记错符号是这套公式最常见的出错点。
\end{yizhu}

(2) 链映射
$
S:
C_{\bullet}(I)\otimes C_{\bullet}(J)
\rightarrow
C_{\bullet}(I\times J)
$
诱导上链复形之间的映射
\begin{align*}
S^*:
C_{\times}^n
&\rightarrow
C_{\otimes}^n
\\
\quad
h
=
\Big(
h_{\alpha_0\cdots\alpha_t}^{i,j}
\Big)
_{\substack{
i,j\in\mathbb{Z}
\\
\alpha_0<\cdots<\alpha_t
}}
&\mapsto
S^*h
:=
\Big(
(S^*h)_{i_0\cdots i_p;j_0\cdots j_r}^{i,j}
\Big)
_{\substack{
i,j\in\mathbb{Z}\phantom{,\ p+r=t}
\\
p,r\in\mathbb{Z},\ p+r=t
\\
i_0<\cdots<i_p\quad\
\\
j_0<\cdots<j_r\quad\
}},
\end{align*}
其分量由
$$
(S^*h)_{i_0\cdots i_p;j_0\cdots j_r}^{i,j}
:=
\sum_{\substack{
\alpha_0\cdots\alpha_t
}}
(-1)^{qr+qs+si}
S_{i_0\cdots i_p;j_0\cdots j_r}^{\alpha_0\cdots\alpha_t}
\boxtimes^{-1}\big(h_{\alpha_0\cdots\alpha_t}^{i,j}\big),
$$
定义，其中系数
$
S_{i_0\cdots i_p;j_0\cdots j_r}^{\alpha_0\cdots\alpha_t}
\in\mathbb{Z}
$
由公式
$$
S\big((i_0,\dots,i_p)\otimes(j_0,\dots,j_r)\big)
=
\sum_{\substack{
\alpha_0<\cdots<\alpha_t
}}
S_{i_0\cdots i_p;j_0\cdots j_r}^{\alpha_0\cdots\alpha_t}
(\alpha_0,\dots,\alpha_t)
$$
确定，并且\newline
\resizebox{0.98\textwidth}{!}{$\displaystyle
\boxtimes:
\Hom_{\VB_G(U_{i_0\cdots i_p})}
\big(
(\mathcal{F}^i,\theta^i)|_{U_{i_0\cdots i_p}},
(\mathcal{F}'^{i+q},\theta'^{i+q})|_{U_{i_0\cdots i_p}}
\big)
\underset{\field}{\otimes}
\Hom_{\VB_H(V_{j_0\cdots j_r})}
\big(
(\mathcal{G}^j,\eta^j)|_{V_{j_0\cdots j_r}},
(\mathcal{G}'^{j+s},\eta'^{j+s})|_{V_{j_0\cdots j_r}}
\big)
$}\\
\resizebox{0.98\textwidth}{!}{$\displaystyle
\quad
\xrightarrow{\cong}
\Hom_{\VB_{G\times H}(U_{i_0\cdots i_p}\times V_{j_0\cdots j_r})}
\Big(
\big(
(\mathcal{F}^i,\theta^i)
\boxtimes
(\mathcal{G}^j,\eta^j)
\big)
\big|_{U_{i_0\cdots i_p}\times V_{j_0\cdots j_r}},
\big(
(\mathcal{F}'^{i+q},\theta'^{i+q})
\boxtimes
(\mathcal{G}'^{j+s},\eta'^{j+s}
\big)
\big|_{U_{i_0\cdots i_p}\times V_{j_0\cdots j_r}}
\Big)
$}\\
就是命题 \ref{prop:BoxTensorFunctorFullyFaithful} 中给出的同构。由 $S$ 的定义
（定义 \ref{defn:EZMaps}.(1)）可知
\begin{center}
\resizebox{0.98\textwidth}{!}{$\displaystyle
(S^*h)_{i_0\cdots i_p;j_0\cdots j_r}^{i,j}
=
(-1)^{qr+qs+si}
\boxtimes^{-1}
\bigg(
\sum_{
0=l_0\le l_1\le\cdots l_{p+1}=r
}
(-1)^{l_1+\cdots+l_p}
h_{
(i_0,j_0),\dots,(i_0,j_{l_0}),
(i_1,j_{l_1}),\dots,(i_1,j_{l_2}),
\dots,
(i_p,j_{l_p}),\dots,(i_p,j_r)
}
^{i,j}
\bigg).
$}
\end{center}

\noindent
(3) 链同伦
$
\mathbb{H}:
C_{\bullet}(I\times J)
\rightarrow
C_{\bullet}(I\times J)
$
诱导上链复形之间的映射
$$
\mathbb{H}^{\ast}\colon C_{\times}^{\bullet}\to C_{\times}^{\bullet -1}
$$其中对每个 $f\coloneqq\left(f_{\beta_0\dotsb\beta_{t+1}^k}\right)_
{
    \begin{subarray}{c}
        k\in\mathbb{Z}\\ 
        \beta_0<\dotsb<\beta_{t+1}
    \end{subarray}
    }$，有$$
(\mathbb{H}^{\ast} f)\coloneqq\left((\mathbb{H}^{\ast} f)^k_{\alpha_0\dotsb\alpha_t}\coloneqq\sum_{\beta_0<\dotsb<\beta_{t+1}}\left. \mathbb{H}_{\alpha_0\dotsb\alpha_t}^{\beta_0\dotsb\beta_{t+1}} f_{\beta_0\dotsb\beta_{t+1}}^k\right|_{W_{\alpha_0\dotsb\alpha_t}}\right)
_{
    \begin{subarray}{c}
        k\in\mathbb{Z}\\ 
        \alpha_0<\dotsb<\alpha_t
    \end{subarray}
}
$$这里记 \[
    \mathbb{H}((\alpha_0,\dotsb,\alpha_t))\eqqcolon\sum_{\beta_0<\dotsb<\beta_{t+1}}\mathbb{H}_{\alpha_0\dotsb\alpha_t}^{\beta_0\dotsb\beta_{t+1}} (\beta_0,\dotsb ,\beta_{t+1}).
\]
\end{dfn}

\begin{yizhu}[原文 (3) 与证明里的几处笔误]
(3) 中 $f\coloneqq\big(f_{\beta_0\dotsb\beta_{t+1}^k}\big)$ 的上标 $k$ 被误排进了下标，应读作
$f_{\beta_0\dotsb\beta_{t+1}}^{k}$；同段公式里的 $f_{\beta_0\dotsb\beta_{t+1}}^k$ 才是正确写法。
正文说「映射 $S$、$T$ 与 $H$」，而图中和定义里写作 $\mathbb{H}$，二者指同一个同伦。最后，
证明末尾引用「命题 \ref{prp:STHd-relation}」是自指，实际用到的是命题
\ref{prp:chain-homotopy-TS} 中 $(S,T,\mathbb{H})$ 的关系。译文均照原文排印。
\end{yizhu}

现在可以验证
\begin{prp}\label{prp:STHd-relation}
    \begin{itemize}
        \item 记 $\delta$ 为 \v{C}ech 微分，则 $\delta\mathbb{H}^{\ast} +\mathbb{H}^{\ast} \delta=S^{\ast} T^{\ast} -\id$。
        \item $S^{\ast}$ 保持复合，即 $S^{\ast} (f\comp g)=S^{\ast} f\comp S^{\ast} g$。
        \item 记 $d_{\mathcal{H}^{\bullet}}$ 为内层 hom 复形的微分，则 $d_{\mathcal{H}^{\bullet}}\mathbb{H}^{\ast} =\mathbb{H}^{\ast}d_{\mathcal{H}^{\bullet}}$。
        \item $dT^{\ast} =T^{\ast} d$。
        \item $T^{\ast} S^{\ast} =\id$。
    \end{itemize}
\end{prp}
\begin{proof}
    这些都是直接验证，用到定义以及命题 \ref{prp:STHd-relation} 中的关系。
\end{proof}

\begin{keytakeaways}
\item 对有序集 $I$、$J$，洗牌映射 $S$ 与 Alexander--Whitney 映射 $T$ 互为同伦逆，同伦
$\mathbb{H}$ 把二者的差异写成 $\mathbb{H}\comp\partial+\partial\comp\mathbb{H}=T\comp S-\id$；
三元组 $(S,T,\mathbb{H})$ 就是 Eilenberg--Zilber 收缩。
\item 把这三张链映射作用到 \v{C}ech 复形上，得到 $C_{\otimes}^{\bullet}$ 与
$C_{\times}^{\bullet}$ 之间的上链映射 $S^{\ast}$、$T^{\ast}$、$\mathbb{H}^{\ast}$；系数
$T_{\alpha_0\cdots\alpha_t}^{i_0\cdots i_p;j_0\cdots j_r}$ 与
$S_{i_0\cdots i_p;j_0\cdots j_r}^{\alpha_0\cdots\alpha_t}$ 完全由链映射在基上的取值决定。
\item 命题 \ref{prp:STHd-relation} 的五条关系说明它在 dg 层面仍是收缩：$T^{\ast}S^{\ast}=\id$，$S^{\ast}T^{\ast}$
与恒等只差 $\delta\mathbb{H}^{\ast}+\mathbb{H}^{\ast}\delta$，$S^{\ast}$ 保持复合，$T^{\ast}$ 与
微分交换，因此两个 Hom 复形拟同构。
\item 公式里反复出现的 $(-1)^{qr+qs+si}$ 都来自分次交换（Koszul 符号）；核对符号时按
「先越过 $r+s$、再越过 $s$」拆开最省事。
\end{keytakeaways}

\begin{conceptbox}{dg 范畴与 $A_\infty$ 范畴（回顾）}
dg 范畴的态射空间是上链复形、合成严格结合；$A_\infty$ 范畴只要求合成在链同伦意义下结合，高阶运算 $\mu^d$ 记录结合子偏离严格结合的程度。两者可通过「同伦转移」互相转换（Keller、Seidel），而附录 A 的 Eilenberg--Zilber 收缩正是构造这种转移的技术工具——这也解释了为什么扭 Künneth 公式的证明要在这里落脚。
\par\smallskip\noindent{\footnotesize 参考：Keller, arXiv:math/0601185；Seidel, \emph{Fukaya categories and Picard--Lefschetz theory}, 2008。}
\end{conceptbox}
